\documentclass[11pt]{article}

\usepackage[margin=1in]{geometry}
\usepackage{amsmath}
\usepackage{amssymb}
\usepackage{mathtools}
\usepackage{enumitem}
\usepackage{parskip}
\usepackage{graphicx}
\usepackage{xcolor}
\usepackage{listings}
\usepackage{hyperref}

\newcommand{\flops}{\text{FLOPs}}
\newcommand{\flopss}{\text{FLOPs/s}}
\newcommand{\Tmath}{T_{\text{math}}}
\newcommand{\Tcomm}{T_{\text{comm}}}
\newcommand{\Tcomp}{T_{\text{comp}}}
\newcommand{\Wici}{W_{\text{ici}}}
\newcommand{\Whbm}{W_{\text{hbm}}}

\definecolor{codegray}{rgb}{0.95,0.95,0.95}
\definecolor{codecomment}{rgb}{0,0.5,0}
\definecolor{codekeyword}{rgb}{0.6,0,0.6}
\definecolor{codestring}{rgb}{0.58,0,0.82}

\lstset{
  basicstyle=\ttfamily\footnotesize,
  backgroundcolor=\color{codegray},
  commentstyle=\color{codecomment},
  keywordstyle=\color{codekeyword},
  stringstyle=\color{codestring},
  language=Python,
  breaklines=true,
  showstringspaces=false,
  frame=single,
  rulecolor=\color{black!30},
  numbers=left,
  numberstyle=\tiny\color{gray},
  numbersep=8pt,
}

\title{Scaling Book --- Section 9 Exercises}
\author{Rahul Bir}
\date{}

\begin{document}
\maketitle

\section*{Q projection op exercise}

\[
X[32, 1024, 8192] \cdot W_Q[8192, 32, 256] \quad \text{on v2 } 4 \times 2 \text{ slice.}
\]

\[
\text{BS} = \frac{32 \cdot 1024}{4} = 8192 \quad \text{so we're compute bound.}
\]

\[
\text{FLOPS} = 2 \cdot 32 \cdot 1024 \cdot 8192 \cdot 32 \cdot 256 \approx 2.46 \cdot 10^{12}.
\]

\[
T_{\text{comp}} = \frac{2.46 \cdot 10^{12}}{8 \cdot 23 \cdot 10^{12}} = 13.4 \text{ ms.}
\]

This is close to the wall duration in the profile!

\section*{Worked Problems}

\section*{Q1) Trace}

\paragraph{multiply\_reduce\_fusion.1:}
\[
\underbrace{\text{bf16}[4096, 8192]}_{W_1} \cdot \underbrace{\text{bf16}[8192]}_{\text{In}} \to \underbrace{\text{bf16}[4096]}_{\text{Tmp}} \quad \Rightarrow \text{matmul}
\]

\paragraph{multiply\_reduce\_fusion:}
\[
\underbrace{\text{bf16}[8192, 4096]}_{W_2} \cdot \underbrace{\text{bf16}[4096]}_{\text{Tmp}} \to \underbrace{\text{bf16}[8192]}_{\text{Out}} \quad \Rightarrow \text{matmul (feeding in result from last HLO op)}
\]

\paragraph{all\_reduce:}
\[
\text{all\_reduce}(\text{bf16}[1, 8192]) \to \text{bf16}[1, 8192]
\]

Running on a v5e-8 slice and replica groups is 8, so sharded across all devices on contracting dim.

\paragraph{Computation:}
\[
W_1[I_{xy}, J] \cdot_J \text{In}[J] \to \text{Tmp}[I_{xy}]
\]
\[
W_2[K, I_{xy}] \cdot_I \text{Tmp}[I_{xy}] \to \text{Out}[K]\{U_{xy}\}
\]
\[
\text{AllReduce}_{xy}(\text{Out}[K]\{U_{xy}\}) \to \text{Out}[K]
\]

\[
\underbrace{\text{bf16}[32768, 8192]}_{W_1} \cdot \underbrace{\text{bf16}[8192]}_{\text{In}} \cdot \underbrace{\text{bf16}[8192, 32768]}_{W_2} \to \text{bf16}[8192]
\]

This operation is essentially the up and down projection in an MLP where the $F$ dim is sharded across all 8 devices.

\section*{Q2) Profiling \& Sharding Optimization}

The following is the full Jupyter notebook used to profile and re-shard the forward pass. Each code block corresponds to one cell in the notebook.

\subsection*{Cell 1: Imports, Config, Layer, Weights}

\begin{lstlisting}[language=Python]
from flax import struct
from typing import Callable, Tuple

import collections

import jax
import jax.numpy as jnp

P = jax.sharding.PartitionSpec

ShardingRules = collections.namedtuple('ShardingRules',
 ['batch', 'sequence', 'd_model', 'query_heads', 'key_heads', 'key_dim', 'd_ff', 'vocab'])

def _logical_to_physical(logical: Tuple[str, ...], rules: ShardingRules):
    """Converts logical to physical pspec."""
    return P(*(getattr(rules, l) for l in logical))

def _logical_to_sharding(logical: Tuple[str, ...], mesh: jax.sharding.Mesh, rules: ShardingRules):
    """Converts logical to sharding."""
    return jax.sharding.NamedSharding(mesh, _logical_to_physical(logical, rules))


@struct.dataclass
class Config:
  d_model: int
  ffw_multiplier: int
  num_layers: int

  query_heads: int
  kv_heads: int
  key_dim: int

  vocab_size: int

  dtype: jnp.dtype = jnp.bfloat16


@struct.dataclass
class TensorInfo:
  shape: jax.ShapeDtypeStruct
  logical_axes: tuple[str, ...]
  initializer: Callable | None = None


@struct.dataclass
class Layer:

  q: jax.Array | TensorInfo
  k: jax.Array | TensorInfo
  v: jax.Array | TensorInfo
  proj: jax.Array | TensorInfo

  w1: jax.Array | TensorInfo
  w2: jax.Array | TensorInfo

  gamma1: jax.Array | TensorInfo
  gamma2: jax.Array | TensorInfo

  @classmethod
  def abstract(cls, cfg: Config):
    return Layer(
        q=TensorInfo(
            jax.ShapeDtypeStruct((cfg.d_model, cfg.query_heads, cfg.key_dim), dtype=cfg.dtype),
            ('d_model', 'query_heads', 'key_dim'),
            jax.nn.initializers.he_normal(in_axis=0, out_axis=(1,2)),
        ),
        k=TensorInfo(
            jax.ShapeDtypeStruct((cfg.d_model, cfg.kv_heads, cfg.key_dim), dtype=cfg.dtype),
            ('d_model', 'query_heads', 'key_dim'),
            jax.nn.initializers.he_normal(in_axis=0, out_axis=(1,2)),
        ),
        v=TensorInfo(
            jax.ShapeDtypeStruct((cfg.d_model, cfg.kv_heads, cfg.key_dim), dtype=cfg.dtype),
            ('d_model', 'query_heads', 'key_dim'),
            jax.nn.initializers.he_normal(in_axis=0, out_axis=(1,2)),
        ),
        proj=TensorInfo(
            jax.ShapeDtypeStruct((cfg.query_heads, cfg.key_dim, cfg.d_model), dtype=cfg.dtype),
            ('query_heads', 'key_dim', 'd_model'),
            jax.nn.initializers.he_normal(in_axis=(0,1), out_axis=2),
        ),
        w1=TensorInfo(
            jax.ShapeDtypeStruct((cfg.d_model, cfg.d_model * cfg.ffw_multiplier), dtype=cfg.dtype),
            ('d_model', 'd_ff'),
            jax.nn.initializers.he_normal(in_axis=0, out_axis=1),
        ),
        w2=TensorInfo(
            jax.ShapeDtypeStruct((cfg.d_model * cfg.ffw_multiplier, cfg.d_model), dtype=cfg.dtype),
            ('d_ff', 'd_model'),
            jax.nn.initializers.he_normal(in_axis=0, out_axis=1),
        ),
        gamma1=TensorInfo(
            jax.ShapeDtypeStruct((cfg.d_model,), dtype=cfg.dtype),
            ('d_model',),
            jax.nn.initializers.constant(1.0),
        ),
        gamma2=TensorInfo(
            jax.ShapeDtypeStruct((cfg.d_model,), dtype=cfg.dtype),
            ('d_model',),
            jax.nn.initializers.constant(1.0),
        ),
    )

@struct.dataclass
class Weights:
  layers: list[Layer]
  embedding: jax.Array | TensorInfo

  @classmethod
  def abstract(cls, cfg: Config):
      return Weights(
        layers=[Layer.abstract(cfg) for _ in range(cfg.num_layers)],
        embedding=TensorInfo(
          jax.ShapeDtypeStruct((cfg.vocab_size, cfg.d_model), cfg.dtype),
          ('vocab', 'd_model'),
          jax.nn.initializers.he_normal(in_axis=0, out_axis=1)
        ),
    )

  @classmethod
  def sharding(cls, cfg: Config, mesh: jax.sharding.Mesh, rules: ShardingRules):
    abstract = cls.abstract(cfg)
    return jax.tree.map(lambda info: _logical_to_sharding(info.logical_axes, mesh, rules), abstract,
                        is_leaf=lambda x: isinstance(x, TensorInfo))

  @classmethod
  def init(cls, cfg: Config, key: jax.random.PRNGKey, mesh: jax.sharding.Mesh, rules: ShardingRules):
    abstract = cls.abstract(cfg)

    def _init():
      num_leaves = len(jax.tree.leaves(abstract, is_leaf=lambda x: isinstance(x, TensorInfo)))
      rng_iter = iter(jax.random.split(key, num_leaves))
      return jax.tree.map(
          lambda info: info.initializer(next(rng_iter), info.shape.shape, info.shape.dtype), abstract,
          is_leaf=lambda x: isinstance(x, TensorInfo))

    sharding = cls.sharding(cfg, mesh, rules)
    return jax.jit(_init, out_shardings=sharding)()
\end{lstlisting}

\subsection*{Cell 2: rms\_norm, layer\_forward, forward}

\begin{lstlisting}[language=Python]
def rms_norm(x: jax.Array, gamma: jax.Array) -> jax.Array:
    """Apply RMS normalization."""
    """
    ORIGINAL:
    Negligible flops but for sharding scheme:
    x[B, T, D@X] => x[B, T] by fused vector ops in the last matmul and then an AR across the partial sums for the mean
    outputs similar sharding x[B, T, D@X]

    INSTEAD:
    Still need to AR but now that B is sharded as well, this comm should be faster.
    """
    rms = jax.lax.rsqrt(jnp.mean(jnp.square(x), axis=-1, keepdims=True) + 1e-6)
    return gamma * x * rms

def layer_forward(x: jax.Array, layer: Layer) -> jax.Array:
  x = jax.lax.with_sharding_constraint(x, jax.P('x', None, 'y'))

  with jax.named_scope('pre_attn_norm'):
    attn_in = rms_norm(x, layer.gamma1)

  # Q, K, V projections
  with jax.named_scope('qkv_matmul'):
    """
    ORIGINAL:
    roofline = 3*2*btdhn / (8 * 1.97e14) = 1ms
    actual = 8.174ms
    Note: (N == Q == K == V)

    op: x[B, T, D@X] @ W_q[D@X, N@Y, H] => q,k,v[N@Y, H, B, T]
    steps:
    x[B, T, D@X] @ W_q[D@X, N@Y, H] => q,k,v[N@Y, H, B, T]{U_X}
    AR_X(q,k,v[N@Y, H, B, T]{U_X}) = q,k,v[N@Y, H, B, T]
    (AR dominates the time here)

    INSTEAD:
    If x[B@X, T, D@Y] (like in the mixed FSDP+TP section):
    x[B@X, T, D@Y] @ W_q[D@X, N@Y, H] => x[B@X, T, N@Y, H]
    (still need to do 2 AG which is still slightly less than AR since
     one AG is over Y while the AR is over X. Additionally, one AG
     can be done ahead of time if optimized by the compiled.)
    """
    q = jnp.einsum('btd,dhq->bhtq', attn_in, layer.q)
    k = jnp.einsum('btd,dhk->bhtk', attn_in, layer.k)
    v = jnp.einsum('btd,dhv->bhtv', attn_in, layer.v)

  # TODO: add RoPE here

  # just to ensure
  x = jax.lax.with_sharding_constraint(x, jax.P('x', None, 'y'))

  # Attention
  with jax.named_scope('attn'):
    """
    ORIGINAL:
    roofline = (2bhtsn + 2bhstn) / (8 * 1.97e14) = 0.08ms
    actual = 1ms

    (already very small, not dominant flops)
    """
    scale = q.shape[-1] ** -0.5
    # TODO: add masking here

    num_query_heads, num_kv_heads = q.shape[1], k.shape[1]

    if num_query_heads == num_kv_heads or num_kv_heads == 1:
      qk = jnp.einsum('bhtd,bhsd->bhts', q, k) * scale
      logits = jax.nn.softmax(qk.astype(jnp.float32), axis=-1)
      attn_vec = jnp.einsum('bhsd,bhts->bhtd', v, logits)
    else:
      assert num_query_heads % num_kv_heads == 0
      q = q.reshape(q.shape[0:1] +
                    (num_kv_heads, num_query_heads // num_kv_heads) +
                    q.shape[2:])
      qk = jnp.einsum('bqhtd,bhsd->bqhts', q, k) * scale
      logits = jax.nn.softmax(qk.astype(jnp.float32), axis=-1)
      attn_vec = jnp.einsum('bqsd,bqhts->bqhtd', v, logits)
      attn_vec = attn_vec.reshape(attn_vec.shape[0:1] + (num_query_heads,) + attn_vec.shape[3:])

  # Attention proj
  """
  ORIGINAL:
  roofline = 2bhtnd / (8 * 1.97e14) = 0.3ms
  actual = 1.98ms
  (dominated by AR)
  op: attn_vec[B, H, T, N@Y] @ [H, N@Y, D@X] => [B, T, D@X]
  steps:
  attn_vec[B, H, T, N@Y] @ [H, N@Y, D@X] => [B, T, D@X]{U_Y}
  AR_Y([B, T, D@X]{U_Y}) => [B, T, D@X]

  INSTEAD:
  op: attn_vec[B@X, H, T, N@Y] @ [H, N@Y, D@X] => [B@X, T, D@Y]
  steps:
  AG_X([H, N@Y, D@X]) => [H, N@Y, D]
  attn_vec[B@X, H, T, N@Y] @ [H, N@Y, D] => [B@X, T, D]{U_Y}
  RS_Y,D([B@X, T, D]{U_Y}) => [B@X, T, D@Y]

  Timing doesn't change that much under sharding.
  """
  with jax.named_scope('attn_proj'):
    attn_out = jnp.einsum('bhtv,hvd->btd', attn_vec, layer.proj)

  # Residual connection
  with jax.named_scope('residual'):
    x = x + attn_out

  # Second RMSNorm (Pre-LN for FFN)
  with jax.named_scope('ffn_pre_norm'):
    ffw_in = rms_norm(x, layer.gamma2)

  jax.lax.with_sharding_constraint(ffw_in, jax.P('x', None, 'y'))

  """
  ORIGINAL:
  roofline = 2*2btdf / (8 * 1.97e14) = 5ms
  actual = 24.9ms (much larger!)

  op1: [B, T, D@X] @ [D@X, F@Y] = [B, T, F@Y]
  steps:
  [B, T, D@X] @ [D@X, F@Y] = [B, T, F@Y]{U_X}
  AR_X([B, T, F@Y]{U_X}) = [B, T, F@Y]

  op2: [B, T, F@Y] @ [F@Y, D@X] => [B, T, D@X]
  steps:
  [B, T, F@Y] @ [F@Y, D@X] => [B, T, D@X]{U_Y}
  AR_Y([B, T, D@X]{U_Y}) = [B, T, D@X]

  INSTEAD:
  op1: [B@X, T, D@Y] @ [D@X, F@Y] => [B@X, T, F@Y]
  steps:
  AG_Y([B@X, T, D@Y]) = [B@X, T, D]
  AG_X([D@X, F@Y]) = [D, F@Y]
  [B@X, T, D] @ [D, F@Y] = [B@X, T, F@Y]

  op2: [B@X, T, F@Y] @ [F@Y, D@X] => [B@X, T, D@Y]
  steps:
  AG_X([F@Y, D@X]) = [F@Y, D]
  [B@X, T, F@Y] @ [F@Y, D] => [B@X, T, D]{U_Y}
  RS_Y,D([B@X, T, D]{U_Y}) = [B@X, T, D@Y]

  (under the new scheme, we would have cheaper comms)
  """
  with jax.named_scope('ffw'):
    ffw_out = jnp.einsum('btd,df->btf', ffw_in, layer.w1).astype(jnp.bfloat16)
    ffw_out = jax.nn.gelu(ffw_out)
    ffw_out = jnp.einsum('btf,fd->btd', ffw_out, layer.w2).astype(jnp.bfloat16)

  # Residual connection
  with jax.named_scope('residual'):
    x = x + ffw_out

  # just to ensure
  x = jax.lax.with_sharding_constraint(x, jax.P('x', None, 'y'))

  return x


def forward(x: jax.Array, weights: Weights) -> jax.Array:
  """Forward pass through the network."""


  # Initially, x = [B@X, T]
  one_hot = jax.nn.one_hot(x, cfg.vocab_size)
  """
  ORIGINAL
  x = [B@X, T, V]
  bottom matmul should take (2btvd) / (8 * 1.97e14) = 2.7ms
  actual time = 11.4ms!

  op: x[B@X, T, V] @ [V, D@X] => [B, T, D@X]
  steps:
  AG_X(x[B@X, T, V]) => x[B, T, V]
  x[B, T, V] @ [V, D@X] => [B, T, D@X]


  INSTEAD
  op: x[B@X, T, V] @ [V, D@X] => [B@X, T, D@Y]
  steps:
  Reshard([V, D@X]) => [V, D@Y]
  x[B@X, T, V] @ [V, D@Y] => [B@X, T, D@Y]
  """
  x = jnp.einsum('vd,btv->btd', weights.embedding, one_hot)

  for idx, layer in enumerate(weights.layers):
    with jax.named_scope(f'layer_{idx}'):
      x = layer_forward(x, layer)

  logits = jnp.einsum('vd,btd->btv', weights.embedding, x)
  return jax.nn.log_softmax(logits, axis=-1)
\end{lstlisting}

\subsection*{Cell 3: Sharding rules}

\begin{lstlisting}[language=Python]
fsdp_rules = ShardingRules(
    batch=('x', 'y'),
    sequence=None,
    d_model=('x', 'y'),
    query_heads=None,
    key_heads=None,
    key_dim=None,
    d_ff=None,
    vocab=None
)

model_parallel_rules = ShardingRules(
    batch=None,
    sequence=None,
    d_model=None,
    query_heads=('x', 'y'),
    key_heads=('x', 'y'),
    key_dim=None,
    d_ff=('x', 'y'),
    vocab=None
)

mixed_rules = ShardingRules(
    batch='x',
    sequence=None,
    d_model='x',
    query_heads='y',
    key_heads='y',
    key_dim=None,
    d_ff='y',
    vocab=None
)
\end{lstlisting}

\subsection*{Cell 4: Clear live arrays}

\begin{lstlisting}[language=Python]
for arr in jax.live_arrays():
  arr.delete()
\end{lstlisting}

\subsection*{Cell 5: Config, mesh, init weights}

\begin{lstlisting}[language=Python]
Auto = jax.sharding.AxisType.Auto

cfg = Config(
    d_model=8192,
    ffw_multiplier=4,
    num_layers=4,
    query_heads=16,
    kv_heads=16,
    key_dim=256,
    vocab_size=32_128,
    dtype=jnp.bfloat16,
)

mesh = jax.make_mesh((4, 2), ('x', 'y'), (Auto, Auto))
jax.set_mesh(mesh)
print(mesh)

rules = mixed_rules

rng = jax.random.PRNGKey(42)
weights = Weights.init(cfg, rng, mesh, rules)
\end{lstlisting}

\subsection*{Cell 6: Compile}

\begin{lstlisting}[language=Python]
batch_size = 8
seq_len = 1024

input_sharding = _logical_to_sharding(('batch', 'sequence'), mesh=mesh, rules=rules)
x = jnp.zeros((batch_size, seq_len), dtype=jnp.int32, device=input_sharding)

compiled = jax.jit(forward).lower(x, weights).compile()
\end{lstlisting}

\subsection*{Cell 7: Profile}

\begin{lstlisting}[language=Python]
with jax.profiler.trace("/kaggle/working/tensorboard"):
  logits = compiled(x, weights)
  jax.block_until_ready(logits)
\end{lstlisting}

\subsection*{Cell 8: Zip the tensorboard trace}

\begin{lstlisting}[language=Python]
import os
import shutil
from IPython.display import FileLink, display

# Folder where JAX wrote the TensorBoard trace
src_dir = "/kaggle/working/tensorboard"

# Zip output path
zip_base = "/kaggle/working/tensorboard"
zip_file = zip_base + ".zip"

# Check trace folder exists
if not os.path.exists(src_dir):
    raise FileNotFoundError(f"{src_dir} does not exist. Make sure your profiler wrote here.")

# Remove old zip if present
if os.path.exists(zip_file):
    os.remove(zip_file)

# Create zip
shutil.make_archive(
    base_name=zip_base,
    format="zip",
    root_dir=src_dir,
)

# Verify it exists
print("Created:", zip_file)
!ls -lh /kaggle/working/tensorboard.zip

# Make Kaggle download link
os.chdir("/kaggle/working")
display(FileLink("tensorboard.zip"))
\end{lstlisting}

\subsection*{Additional Questions}

Under this fix, the highest MXU I now get is 12.7\% compared to 13.3\%. However, the total forward pass time has halved from 206 ms to 109 ms. I suspect the slight decrease in MXU might be from more comms launched even though each comm is much cheaper.

\begin{enumerate}[leftmargin=*]
  \item The sharding strategy I changed it to is a mixed FSDP + TP. Specifically, this is 4-way FSDP and 2-way TP.
  \item $\text{BS} = 8$ but per device it's $\text{BS}_{\text{local}} = 2$. $d_{\text{model}} = 8192$ but per device (on activations) becomes 2048. $d_{\text{ff}} = 4 \cdot d_{\text{model}} = 32768$ but per device is 16384.
  \item Including the projections in attention, under my optimized scheme, the attention block takes around 7 ms while the MLP block takes around 11 ms.
  \item Rooflines are written above their corresponding code blocks above (BS and sequence length are large enough that on roofline we should be compute bound).
\end{enumerate}

\end{document}
